\documentclass[11pt]{amsart}
\usepackage{amsmath,amssymb,amsthm,mathtools}
\usepackage[margin=1.12in]{geometry}
\usepackage{hyperref}
\hypersetup{hidelinks}
\usepackage{booktabs}
\usepackage{array}
\usepackage{xcolor}
\usepackage[final,expansion=false]{microtype}
\setlength{\parskip}{4pt plus 1pt}
\emergencystretch=2.5em

\newtheorem{theorem}{Theorem}[section]
\newtheorem{proposition}[theorem]{Proposition}
\newtheorem{lemma}[theorem]{Lemma}
\newtheorem{corollary}[theorem]{Corollary}
\theoremstyle{definition}
\newtheorem{definition}[theorem]{Definition}
\newtheorem{problem}[theorem]{Open closure problem}
\theoremstyle{remark}
\newtheorem{remark}[theorem]{Remark}
\newtheorem{example}[theorem]{Example}

\newcommand{\RH}{\mathrm{RH}}
\newcommand{\Lam}{\Lambda}
\newcommand{\eps}{\varepsilon}
\newcommand{\C}{\mathbb C}
\newcommand{\D}{\mathbb D}
\newcommand{\N}{\mathbb N}
\renewcommand{\Re}{\operatorname{Re}}
\renewcommand{\Im}{\operatorname{Im}}
\DeclareMathOperator{\Res}{Res}
\DeclareMathOperator{\rad}{rad}

\title[The Li sequence as a sampled control system]
{The Li Sequence as a Sampled Control System:\\
FIR Window Stability, Off-Line Modes, and the Arithmetic Closure Problem}

\author{David Alejandro Trejo Pizzo}
\thanks{Independent Researcher. \texttt{dtrejopizzo@gmail.com}}

\begin{document}

\begin{abstract}
We formulate the Li--Keiper sequence of the Riemann zeta function as the
sampled output of a discrete-time spectral system. Under the Cayley map
\[
 w_\rho=1-\rho^{-1},
\]
a zero on the critical line becomes a unit-circle mode, while a zero to the
right becomes an unstable mode \(w_\rho^{-n}\). The symmetric zero orbit has
the exact response
\[
 4-4\cosh(na_\rho)\cos(n\theta_\rho).
\]
Permanent cancellation of an unstable mode by the remaining zeros is
impossible: the exponential rate of the full Li sequence is the reciprocal
of the smallest modulus of an interior Cayley pole.

For every fixed moving-window length \(L\), the output
\[
 Y_L(N)=\sum_{j=0}^{L-1}\lambda_{N+j}
\]
is subexponential if and only if RH holds. More strongly, RH is equivalent
to eventual nonnegativity of every sliding output \(Y_L(N)\). Repeated
window averaging is a cascade of FIR filters
\(H_L(z)=1+\cdots+z^{L-1}\). Since all zeros of \(H_L\) lie on the unit
circle, no such cascade can cancel an off-line mode. We also construct an
adjacent-window bank with no blind spot:
\[
 |H_L(z)|^2+|H_{L+1}(z)|^2\ge \tfrac12|z|^{2L}.
\]

The signal representation lifts exactly to the ordinary von Mangoldt
weights. A sliding window collapses the prime--Laguerre kernel to
\[
 L_{N+L-2}^{(2)}(x)-L_{N-2}^{(2)}(x).
\]
Nested linear averages remain FIR filters and therefore detect, but do not
bound, every unstable mode. A rational off-line quartet proves that a
signed block average may be negative while a pointwise exponential
violation occurs inside the same block. For the deterministic blocks

\[
 I_L=\{L^2,\ldots,L^2+L-1\},
\]

we prove that the \(L^2\)-th root of
\(1+\sum_{n\in I_L}|\lambda_n|\) converges exactly to the largest Cayley
modulus of a zeta zero. The remaining arithmetic task is a
one-sided dissipativity estimate for the complete prime--pole output, or a
nonlinear positive-part estimate. We state that estimate exactly. It is not
proved here; accordingly this paper does not claim a proof of RH. Its
contribution is a rigorous signals-and-systems architecture, exact RH
criteria, and a precise identification of the plant-specific inequality
still missing for the literal weights \(\Lambda(m)\). A lossless prefix
compression reduces the absolute block problem from all \(2^L\) sign
selectors to only \(L\) differences of two Laguerre boundary degrees.
\end{abstract}
\maketitle

\tableofcontents
\newpage

\section{Introduction}

The Riemann Hypothesis is usually presented as a geometric statement about
the zero set of \(\zeta(s)\): every nontrivial zero should lie on
\(\Re s=1/2\). The equivalent Li criterion replaces this infinite geometry
by the sign of a sampled sequence \((\lambda_n)_{n\ge1}\). That replacement
suggests a third language, familiar from electrical engineering.

\begin{center}
\begin{tabular}{>{\raggedright\arraybackslash}p{0.35\textwidth}
                >{\raggedright\arraybackslash}p{0.50\textwidth}}
\toprule
Discrete-time signals and control & Zeta/Li object\\
\midrule
sample number & Li degree \(n\)\\
complex modal frequency & \(w_\rho=1-\rho^{-1}\)\\
unit-circle mode & zero with \(\Re\rho=1/2\)\\
unstable mode & zero with \(\Re\rho>1/2\)\\
sampled plant output & \(\lambda_n\)\\
FIR moving window & \(\sum_{j=0}^{L-1}\lambda_{N+j}\)\\
spectral radius test & \(\limsup|\lambda_n|^{1/n}\)\\
one-sided dissipativity & eventual positivity of a window output\\
\bottomrule
\end{tabular}
\end{center}

The fundamental point is the exact identity
\[
 \left|1-\frac1\rho\right|^2
 =1+\frac{1-2\Re\rho}{|\rho|^2}.
\]
The critical line maps to the unit circle. A zero to the right maps inside
the circle, and the coefficient sequence sees its reciprocal, an unstable
mode. In this language RH asks whether the zeta plant is marginally stable:
polynomial growth is allowed, but a fixed exponential rate is not.

The formulation separates two tasks. The zero-side task is to prove that an
off-line zero creates an unstable output and cannot be hidden forever by
modal interference. We complete it in Sections~\ref{sec:modes}
and~\ref{sec:global-modes}. The
arithmetic-side task is to prove, from the literal weights \(\Lambda(m)\),
that a suitable filtered output is subexponential or one-sided dissipative.
We derive its exact kernel in Section~\ref{sec:arithmetic}. That second
estimate remains open.

This distinction is essential. Computation of finitely many critical zeros
tests the response of known unit-circle modes, but does not control an
unknown unstable mode above the verified height. Conversely, an FIR filter
can be proved to detect every unstable mode without thereby proving that the
actual plant has none. Detection and exclusion are separate steps.

\section{Li coefficients and the Cayley plant}
\label{sec:modes}

Let
\[
 \xi(s)=\frac12s(s-1)\pi^{-s/2}\Gamma(s/2)\zeta(s)
\]
be the completed Riemann function. The Li coefficients are given by
\[
 \lambda_n=\sum_\rho\left\{1-\left(1-\frac1\rho\right)^n\right\},
                                                        \tag{2.1}
\]
with the standard symmetric interpretation of the zero sum. Equivalently,
\[
 G(z):=z\frac{d}{dz}\log\xi\!\left(\frac1{1-z}\right)
      =\sum_{n\ge1}\lambda_nz^n.                         \tag{2.2}
\]

\begin{lemma}[Cayley modulus]
For \(\rho=\beta+i\gamma\) and \(w_\rho=1-\rho^{-1}\),
\[
 |w_\rho|^2=1+\frac{1-2\beta}{|\rho|^2}.                 \tag{2.3}
\]
Thus \(|w_\rho|=1\) exactly on the critical line and
\(|w_\rho|<1\) exactly when \(\beta>1/2\).
\end{lemma}

\begin{proof}
Direct expansion gives
\[
 |w_\rho|^2
 =\frac{(\beta-1)^2+\gamma^2}{\beta^2+\gamma^2}
 =1+\frac{1-2\beta}{\beta^2+\gamma^2}.
\]
\end{proof}

\section{FIR windows and filtered Li outputs}

For a finite real impulse response
\[
 h=(h_0,\ldots,h_{L-1}),
\]
define
\[
 H_h(z)=\sum_{j=0}^{L-1}h_jz^j,\qquad
 Y_h(N)=\sum_{j=0}^{L-1}h_j\lambda_{N+j}.                \tag{3.1}
\]
This is a forward-index convention. A mode \(w^{-n}\) is multiplied by
\(H_h(w^{-1})\).

\begin{theorem}[No-cancellation FIR criterion]
\label{thm:fir-general}
Assume that \(H_h(z)\ne0\) for every \(|z|>1\). Then
\[
 \RH\quad\Longleftrightarrow\quad
 \limsup_{N\to\infty}|Y_h(N)|^{1/N}\le1.                 \tag{3.2}
\]
\end{theorem}

\begin{proof}
Under RH, \(Y_h(N)=O_h(N\log N)\). If RH is false, insert
\((2.10)\) into \((3.1)\):
\[
 Y_h(N)
 =-\sum_{|w|=r_0}m_ww^{-N}H_h(w^{-1})+O_h(r_1^{-N}).
                                                        \tag{3.3}
\]
Every multiplier in the dominant sum is nonzero. After equal phases are
grouped, the same mean-square argument as in Theorem~\ref{thm:rate} shows
that the dominant trigonometric polynomial is nonzero. Thus the left side
of \((3.2)\) equals \(r_0^{-1}>1\).
\end{proof}

\subsection{Moving averages}

The unnormalised moving-window filter is
\[
 H_L(z)=1+z+\cdots+z^{L-1}=\frac{1-z^L}{1-z}.            \tag{3.4}
\]
Its zeros are the nontrivial \(L\)-th roots of unity, so it has no zero
outside the unit circle.

\begin{corollary}[Sliding-window stability criterion]
\label{cor:window-stability}
For each fixed \(L\ge1\), put
\[
 Y_L(N)=\sum_{j=0}^{L-1}\lambda_{N+j}.                   \tag{3.5}
\]
Then
\[
 \boxed{\RH\quad\Longleftrightarrow\quad
 \limsup_{N\to\infty}|Y_L(N)|^{1/N}\le1.}                \tag{3.6}
\]
\end{corollary}

This criterion does not require every sample \(\lambda_n\) to have a fixed
sign. It asks only that one FIR output have no unstable exponential rate.

\subsection{Cascades and averages of averages}

Suppose a window output is averaged again, with lengths
\(L_1,\ldots,L_k\). The resulting transfer function is
\[
 H_{\boldsymbol L}(z)=\prod_{\nu=1}^kH_{L_\nu}(z).       \tag{3.7}
\]
All its zeros remain on the unit circle.

\begin{corollary}[Cascade invariance]
\label{cor:cascade}
Every finite cascade of moving averages satisfies
\[
 \RH\quad\Longleftrightarrow\quad
 \limsup_{N\to\infty}|Y_{\boldsymbol L}(N)|^{1/N}\le1.
                                                        \tag{3.8}
\]
No finite number of linear window averages can stabilize or erase an
off-line mode.
\end{corollary}

\begin{remark}
This proves the useful part of averaging windows and then averaging those
averages. The cascade remains a valid detector. It does not prove that the
detector output is subexponential for the actual prime weights. A filter
transfer identity is not a plant estimate.
\end{remark}

\section{One-sided window dissipativity}

Subexponential stability is bilateral. A one-sided sign condition is
sometimes more compatible with arithmetic positivity.

\begin{theorem}[Eventual moving-window positivity]
\label{thm:window-positive}
For every fixed \(L\ge1\), the following are equivalent:
\begin{enumerate}
\item RH;
\item there exists \(N_0\) such that
      \[
        Y_L(N)=\sum_{j=0}^{L-1}\lambda_{N+j}\ge0
        \qquad(N\ge N_0).                                \tag{4.1}
      \]
\end{enumerate}
\end{theorem}

\begin{proof}
Under RH, Li's criterion gives \(\lambda_n\ge0\) for every \(n\), hence
\((4.1)\).

If RH is false, formula \((3.3)\) with \(H_h=H_L\) has a nonzero dominant
real trigonometric polynomial and exponential factor \(r_0^{-N}\). Every
frequency in the polynomial is nontrivial: a nontrivial zeta zero has
nonzero ordinate, so its Cayley phase is not zero modulo \(2\pi\). On the
compact orbit closure the dominant polynomial has Haar mean zero. Since it
is not identically zero, it takes both signs. Minimality on the orbit
closure supplies infinitely many, in fact syndetically many, returns to
each strict sign region. Consequently \(Y_L(N)\) has negative exponential
excursions and cannot satisfy \((4.1)\).
\end{proof}

\begin{remark}
The quantifier for every sufficiently large sliding position is essential.
The existence of some positive long windows does not exclude an off-line
mode: the same unstable sinusoid also has positive phases.
\end{remark}

\section{An adjacent-window filter bank with no blind spot}

A single moving average may annihilate a unit-circle frequency. That is not
dangerous for detecting off-line modes, but it obscures the total
oscillatory content. Two adjacent window lengths remove this blind spot.

\begin{proposition}[Adjacent-bank coercivity]
\label{prop:bank}
For every \(z\in\C\) and \(L\ge1\),
\[
 H_{L+1}(z)-H_L(z)=z^L,                                  \tag{5.1}
\]
and
\[
 \boxed{
 |H_L(z)|^2+|H_{L+1}(z)|^2
 \ge\frac12|z|^{2L}.}                                    \tag{5.2}
\]
\end{proposition}

\begin{proof}
Identity \((5.1)\) is immediate. Since
\[
 |b-a|^2\le2(|a|^2+|b|^2),
\]
take \(a=H_L(z)\), \(b=H_{L+1}(z)\) and use \((5.1)\).
\end{proof}

Define the two-channel energy
\[
 \mathcal E_L(N)=|Y_L(N)|^2+|Y_{L+1}(N)|^2.              \tag{5.3}
\]

\begin{corollary}[Two-channel energy criterion]
\label{cor:energy}
For every fixed \(L\ge1\),
\[
 \RH\quad\Longleftrightarrow\quad
 \limsup_{N\to\infty}\mathcal E_L(N)^{1/(2N)}\le1.
                                                        \tag{5.4}
\]
\end{corollary}

\begin{proof}
Under RH both channels have polynomial growth. If RH is false,
Theorem~\ref{thm:rate} and \((5.2)\) show that the dominant off-line modal
vector cannot vanish in both channels. Its exponential rate is
\(r_0^{-1}>1\).
\end{proof}

Inequality \((5.2)\) is a genuine observability inequality. It solves the
phase-blindness problem of a single window bank. It does not solve the
arithmetic upper-bound problem for \(\mathcal E_L(N)\).

\section{Global modal analysis}
\label{sec:global-modes}

The functional equation and conjugation complete a noncritical zero to
\[
 \rho,\quad\overline\rho,\quad1-\rho,\quad1-\overline\rho.
                                                        \tag{2.4}
\]
Their Cayley images are
\[
 w,\quad\overline w,\quad\overline w^{-1},\quad w^{-1}.
                                                        \tag{2.5}
\]

\begin{theorem}[Exact quartet response]
\label{thm:quartet}
Suppose \(\Re\rho>1/2\), and write
\[
 w_\rho=e^{-a+i\theta},\qquad a>0.
\]
The contribution of the orbit \((2.4)\) to \(\lambda_n\) is
\[
 Q_n(\rho)=4-4\cosh(na)\cos(n\theta).                    \tag{2.6}
\]
It has exponentially large excursions. On the critical line the orbit
coalesces into a conjugate pair whose contribution is
\[
 2-2\cos(n\theta)\in[0,4].                               \tag{2.7}
\]
\end{theorem}

\begin{proof}
Using \((2.5)\),
\[
\begin{aligned}
 Q_n(\rho)
 &=4-\{w^n+\overline w^n+w^{-n}+\overline w^{-n}\}\\
 &=4-2(e^{-na}+e^{na})\cos(n\theta),
\end{aligned}
\]
which is \((2.6)\). Returns of \(n\theta\) to a fixed neighbourhood of zero
give \(\cos(n\theta)\ge1/2\) infinitely often, and then
\[
 Q_n(\rho)\le4-e^{na}.
\]
For irrational rotation these returns are syndetic; for rational rotation
they are periodic. If \(a=0\), the formal quartet is twice the actual
conjugate pair, giving \((2.7)\).
\end{proof}

\begin{theorem}[Global exponential rate]
\label{thm:rate}
If RH is false, define
\[
 r_0=\min_{\Re\rho>1/2}\left|1-\frac1\rho\right|<1.
                                                        \tag{2.8}
\]
Then
\[
 \limsup_{n\to\infty}|\lambda_n|^{1/n}=r_0^{-1}>1.       \tag{2.9}
\]
More precisely, for some \(r_1\in(r_0,1)\),
\[
 \lambda_n=-r_0^{-n}P(n)+O(r_1^{-n}),                   \tag{2.10}
\]
where \(P\) is a nonzero finite real trigonometric polynomial.
\end{theorem}

\begin{proof}
The substitution \(s=(1-z)^{-1}\) maps \(\mathbb D\) onto
\(\Re s>1/2\). A zero \(\rho\) of multiplicity \(m_\rho\) produces in
\((2.2)\) a simple pole at \(w_\rho\), with
\[
 \Res_{z=w_\rho}G(z)=m_\rho w_\rho\ne0.                  \tag{2.11}
\]
The minimum in \((2.8)\) exists because \(|w_\rho|\to1\) as
\(|\Im\rho|\to\infty\). Therefore the radius of convergence of \((2.2)\)
is \(r_0\), and Cauchy--Hadamard proves \((2.9)\).

There are finitely many poles on \(|z|=r_0\). Subtract their principal
parts and continue to a circle of radius \(r_1>r_0\) before the next pole.
Coefficient extraction gives
\[
 \lambda_n=-\sum_{|w|=r_0}m_ww^{-n}+O(r_1^{-n}),
\]
which is \((2.10)\). Grouping equal phases, the Cesaro mean square of
\(P(n)\) equals the sum of the squares of its nonzero grouped coefficients.
Hence \(P\) is bounded away from zero on an infinite subsequence.
\end{proof}

\begin{corollary}[Two-sided syndetic excursions]
\label{cor:two-sided}
If RH is false, there exist \(c>0\), \(R>1\), and two syndetic sets
\(D_+,D_-\subset\mathbb N\) such that, for every sufficiently large index
in the respective set,
\[
 \lambda_n\ge cR^n\quad(n\in D_+),\qquad
 \lambda_n\le-cR^n\quad(n\in D_-).                       \tag{2.12a}
\]
\end{corollary}

\begin{proof}
In \((2.10)\), restrict the dominant real trigonometric polynomial to the
compact closure of its phase orbit. Its Haar mean is zero and it is not
identically zero, so it has two nonempty strict sign regions. Return times
to either open region are syndetic by minimality of the orbit rotation.
The smaller-radius remainder is absorbed after increasing the starting
index.
\end{proof}

\begin{corollary}[Marginal-stability criterion]
\label{cor:raw-stability}
\[
 \RH\quad\Longleftrightarrow\quad
 \limsup_{n\to\infty}|\lambda_n|^{1/n}\le1.              \tag{2.12}
\]
\end{corollary}

\begin{proof}
The reverse implication is Theorem~\ref{thm:rate}. Under RH, the standard
Li asymptotic gives \(\lambda_n=O(n\log n)\).
\end{proof}

\section{Exact lift to the ordinary prime weights}
\label{sec:arithmetic}

We now express the FIR output before passing to zeros. Let
\[
 A_n
 =1-\frac n2\{\gamma+\log(4\pi)\}
  +\sum_{\substack{\ell\ge1\\ \ell\ {\rm odd}}}
   \left\{\left(1-\frac1\ell\right)^n-1+\frac n\ell\right\}
                                                        \tag{6.1}
\]
be the archimedean Li block. For \(\eps>0\), define
\[
\begin{aligned}
 p_n(\eps)
 &=n\sum_{j=1}^n\binom{n-1}{j-1}
       \frac{(-1)^{j-1}}{j\eps^j},\\
 Q_{n,\eps}
 &=\sum_{m\ge2}\frac{\Lambda(m)}{m^{1+\eps}}
       L_{n-1}^{(1)}(\log m),\\
 \lambda_{n,\eps}
 &=A_n+p_n(\eps)-Q_{n,\eps}.
\end{aligned}                                            \tag{6.2}
\]
The complete prime--pole cancellation gives
\[
 \lambda_n=\lim_{\eps\downarrow0}\lambda_{n,\eps}.
                                                        \tag{6.3}
\]

\begin{lemma}[Laguerre window collapse]
\label{lem:lag-window}
For \(N\ge2\) and \(L\ge1\),
\[
 \sum_{j=0}^{L-1}L_{N+j-1}^{(1)}(x)
 =L_{N+L-2}^{(2)}(x)-L_{N-2}^{(2)}(x).                  \tag{6.4}
\]
\end{lemma}

\begin{proof}
Use
\[
 \sum_{k=0}^{M}L_k^{(\alpha)}(x)=L_M^{(\alpha+1)}(x)
\]
at \(M=N+L-2\) and \(M=N-2\), and subtract.
\end{proof}

Set
\[
\begin{aligned}
 A_{N,L}&=\sum_{j=0}^{L-1}A_{N+j},\\
 p_{N,L}(\eps)&=\sum_{j=0}^{L-1}p_{N+j}(\eps),\\
 K_{N,L}(x)&=L_{N+L-2}^{(2)}(x)-L_{N-2}^{(2)}(x).
\end{aligned}                                            \tag{6.5}
\]

\begin{theorem}[Prime--Laguerre FIR identity]
\label{thm:prime-fir}
For every fixed \(N\ge2\) and \(L\ge1\),
\[
\boxed{
 Y_L(N)=\lim_{\eps\downarrow0}
 \left[
 A_{N,L}+p_{N,L}(\eps)
 -\sum_{m\ge2}\frac{\Lambda(m)}{m^{1+\eps}}
        K_{N,L}(\log m)
 \right].}                                              \tag{6.6}
\]
All primes, prime powers, the pole and the archimedean block remain coupled
until after the FIR operation.
\end{theorem}

\begin{proof}
Sum \((6.2)\) over \(n=N,\ldots,N+L-1\), use
Lemma~\ref{lem:lag-window}, and then apply \((6.3)\).
\end{proof}

\begin{corollary}[Exact arithmetic window criterion]
\label{cor:arithmetic-window}
For every fixed \(L\ge1\), RH is equivalent to the existence of \(N_0\)
such that, for all \(N\ge N_0\),
\[
 \lim_{\eps\downarrow0}
 \left[
 A_{N,L}+p_{N,L}(\eps)
 -\sum_{m\ge2}\frac{\Lambda(m)}{m^{1+\eps}}
        K_{N,L}(\log m)
 \right]\ge0.                                           \tag{6.7}
\]
\end{corollary}

\section{Why repeated linear averaging does not close the estimate}

The moving average improves phase variation and retains all unstable modes.
These are useful facts, but neither determines the sign of \((6.6)\).

\subsection{A rational off-line control}

Let
\[
 w=\frac{4i}{5},\qquad
 \rho=\frac1{1-w}=\frac{25+20i}{41}.                    \tag{7.1}
\]
Then \(\Re\rho>1/2\) and \(|w|=4/5\). Its quartet contribution is
\[
 q_n=4-2\Re(w^n+w^{-n})
 =\begin{cases}
 4-2\{(5/4)^n+(4/5)^n\},&n\equiv0\pmod4,\\
 4+2\{(5/4)^n+(4/5)^n\},&n\equiv2\pmod4,\\
 4,&n\ {\rm odd}.
 \end{cases}                                            \tag{7.2}
\]
Every four consecutive sufficiently late samples contain a positive
exponential violation. Nevertheless, for blocks aligned at \(4K+1\),
\[
 \sum_{n=4K+1}^{4K+4M}q_n\longrightarrow-\infty
 \qquad(K\to\infty)                                     \tag{7.3}
\]
for each fixed \(M\ge1\). Thus a signed mean can have the desired upper
sign while hiding a bad sample in every four positions.

\begin{proposition}[Linear-cascade limitation]
\label{prop:linear-limit}
No inference from the sign or boundedness of a signed linear block mean to
a pointwise block barrier is valid on the class of symmetric zero orbits.
Repeating the average does not repair the inference: it only replaces
\(H_L\) by a product of moving-average polynomials, and the off-line mode
survives by Corollary~\ref{cor:cascade}.
\end{proposition}

The control \((7.1)\) is not claimed to arise from the ordinary Euler
product. It falsifies a signal-processing inference, not the desired
arithmetic inequality.

\subsection{Energy also needs an upper bound}

Squaring a filter-bank output prevents sign cancellation:
\(\mathcal E_L(N)\ge0\). But nonnegativity is the wrong orientation. To
exclude an unstable mode one needs an upper bound
\[
 \mathcal E_L(N)\le e^{o(N)}.                            \tag{7.4}
\]
After insertion of \((6.6)\), this square contains separately divergent
prime and polar channels. Expanding them before the Abel limit destroys the
cancellation that defines \(\lambda_n\). Parseval gives an exact detector,
but not estimate \((7.4)\).

\section{Robust nonlinear statistics of the modal growth rate}
\label{sec:robust}

The raw values \(\lambda_n\) are signed and may differ by exponential
factors, so robust estimators should not be applied to them directly. The
scale-invariant sample is
\[
 a_n:=\frac1n\log(1+|\lambda_n|)\ge0.                    \tag{8.1}
\]
Critical behaviour means \(a_n\to0\). An off-line mode gives
\(a_n\ge c>0\) on a syndetic set.

\subsection{Why fixed trimming and the median are blind}

The syndetic excursion set in Corollary~\ref{cor:two-sided} has positive
density, but there is no universal positive lower bound for that density
over all possible phase systems. A median discards every exceptional set
of density below \(1/2\). A trimmed mean that removes a fixed fraction
\(\tau>0\) of the largest observations can likewise remove an off-line
excursion set whose density is smaller than \(\tau\).

The correct robust operation is a vanishing trim. Let
\[
 I_L=\{L^2,L^2+1,\ldots,L^2+L-1\},                       \tag{8.2}
\]
and write the corresponding samples in increasing order,
\[
 a_{L,(1)}\le\cdots\le a_{L,(L)}.
\]
Choose integers \(k_L=o(L)\), and define
\[
\begin{aligned}
 \operatorname{TM}_L
   &=\frac1{L-k_L}\sum_{j=1}^{L-k_L}a_{L,(j)},\\
 \operatorname{Q}_L
   &=a_{L,(L-k_L)}.
\end{aligned}                                            \tag{8.3}
\]

\begin{theorem}[Vanishing-trim criterion]
\label{thm:trim}
For every sequence \(k_L=o(L)\),
\[
 \boxed{
 \RH\Longleftrightarrow
 \operatorname{TM}_L\longrightarrow0
 \Longleftrightarrow
 \operatorname{Q}_L\longrightarrow0.}                   \tag{8.4}
\]
\end{theorem}

\begin{proof}
Under RH, \(\lambda_n=O(n\log n)\), hence
\[
 \max_{n\in I_L}a_n\longrightarrow0.
\]
Both statistics therefore tend to zero.

If RH is false, Corollary~\ref{cor:two-sided} supplies \(c>0\), \(R>1\)
and a syndetic set \(D\) on which \(|\lambda_n|\ge cR^n\). If \(G\) is a
gap bound for \(D\), then every sufficiently large interval of length
\(L\) contains at least \(L/G-O(1)\) members of \(D\). On those members,
\[
 a_n\ge\frac12\log R
\]
for all sufficiently large \(n\). Since \(k_L=o(L)\), deleting the
\(k_L\) largest samples leaves at least \(L/G-o(L)\) such observations.
Consequently
\[
 \liminf_L\operatorname{Q}_L\ge\frac12\log R,\qquad
 \liminf_L\operatorname{TM}_L\ge\frac1{2G}\log R>0.
\]
\end{proof}

\subsection{A geometric-mean criterion}

Define the geometric mean of the per-sample growth factors by
\[
 \operatorname{GM}_L
 =\left\{\prod_{n\in I_L}(1+|\lambda_n|)^{1/n}\right\}^{1/L}.
                                                               \tag{8.5}
\]
Its logarithm is \(L^{-1}\sum_{n\in I_L}a_n\).

\begin{corollary}[Geometric growth criterion]
\label{cor:geom}
\[
 \boxed{\RH\Longleftrightarrow\operatorname{GM}_L\to1.} \tag{8.6}
\]
\end{corollary}

\begin{proof}
Under RH the maximum of the \(a_n\) on \(I_L\) tends to zero. Under
non-RH, the syndetic set used in Theorem~\ref{thm:trim} contributes a
positive lower density of samples with \(a_n\ge\frac12\log R\). Thus
\(\liminf\log\operatorname{GM}_L>0\).
\end{proof}

\subsection{Exact block spectral radius}

The robust statistic above has a sharper unlogged form. Put

\[
 S_L=\sum_{n=L^2}^{L^2+L-1}|\lambda_n|,
 \qquad
 \mathcal R_\zeta
 =\max\left\{1,\max_\rho\left|1-\frac1\rho\right|\right\}.
                                                               \tag{8.7}
\]

If an off-line zero exists, the maximum is attained because the Cayley
moduli tend to one with the zero height.

\begin{theorem}[Block spectral-radius identity]
\label{thm:block-radius}
\[
 \boxed{
 \lim_{L\to\infty}(1+S_L)^{1/L^2}=\mathcal R_\zeta.}    \tag{8.8}
\]
In particular,
\[
 \boxed{
 \RH\quad\Longleftrightarrow\quad
 \log(1+S_L)=o(L^2).}                                   \tag{8.9}
\]
\end{theorem}

\begin{proof}
Under RH, the estimate in the proof of
Corollary~\ref{cor:raw-stability} can be obtained directly by splitting the
critical zeros at height \(n\):
\[
 \lambda_n
 \le4N(n)+n^2\sum_{\gamma>n}\gamma^{-2}
 \ll n\log(n+2).
\]
Thus \(S_L\ll L^3\log(L+2)\), which gives the value one in (8.8).

Suppose RH is false. Use the notation of Theorem~\ref{thm:rate}, so that
\[
 \lambda_n=-R^nP(n)+O(R_1^n),
 \qquad R=r_0^{-1}>R_1,
\]
where, after equal phases have been grouped,
\[
 P(n)=\sum_{j=1}^dM_ju_j^n,\qquad |u_j|=1,\quad u_j\ne u_k.
\]
Uniformly in the initial index \(N\), finite geometric summation gives
\[
 \sum_{n=N}^{N+L-1}|P(n)|^2
 =L\sum_jM_j^2+O_P(1).                                  \tag{8.10}
\]
Hence every sufficiently long block contains an \(n\) for which
\(|P(n)|\ge c_P>0\). Taking \(N=L^2\), the smaller exponential in the
remainder is absorbed and
\[
 cR^{L^2}\le S_L\le CL R^{L^2+L}.                       \tag{8.11}
\]
It follows that \(L^{-2}\log S_L\to\log R\). Functional symmetry pairs
the innermost Cayley pole with the largest reciprocal modulus, so
\(R=\mathcal R_\zeta\). This proves (8.8)--(8.9).
\end{proof}

\begin{remark}
Winsorisation with a fixed clipping fraction has the same blindness as
fixed trimming. A harmonic mean is dominated by samples near zero and can
miss an unstable oscillatory mode. Poisson, negative-binomial and Bayesian
count models require stochastic assumptions not supplied by the zeta
problem. The deterministic compact-group return theorem, rather than an
independence model, is the relevant substitute.
\end{remark}

\section{The nonlinear closure problem}

There are two exact versions of the missing estimate.

\subsection{One-sided dissipativity for a fixed FIR output}

The cleanest signals-and-systems target is
\[
 Y_L(N)\ge0\quad(N\ge N_0).                              \tag{9.1}
\]
It is weaker than coefficientwise Li positivity because individual
\(\lambda_n\) may be negative and be compensated inside the moving window.
It is nevertheless equivalent to RH because the FIR transfer has no zero
off the unit circle.

\subsection{Positive-part supply on growing blocks}

For a growing interval \(I=[N,N+M-1]\), define
\[
 \mathcal S(I)=\sum_{n\in I}(\lambda_n-e^{\sqrt n})_+.
                                                        \tag{9.2}
\]
The pointwise upper barrier holds on \(I\) exactly when
\(\mathcal S(I)=0\). Its arithmetic representation is
\[
\begin{aligned}
\mathcal S(I)
=\lim_{\eps\downarrow0}\sum_{n\in I}
\bigg(
 A_n+p_n(\eps)
 -\sum_{m\ge2}\frac{\Lambda(m)}{m^{1+\eps}}
       L_{n-1}^{(1)}(\log m)
 -e^{\sqrt n}
\bigg)_+ .
\end{aligned}                                            \tag{9.3}
\]
The positive part is taken only after pole, Gamma, primes and all prime
powers have been combined.

\subsection{Absolute-block selector}

The block norm in Theorem~\ref{thm:block-radius} also has a literal-prime
form which keeps the two divergent channels coupled. Define
\[
 d\nu_\eps(u)
 =e^{-\eps u}\,du
 -\sum_{m\ge2}\frac{\Lambda(m)}{m^{1+\eps}}
       \delta_{\log m}(du),                              \tag{9.4}
\]
and, for \(\sigma\in[-1,1]^L\),
\[
 P_{\sigma,L}(u)
 =\sum_{j=0}^{L-1}\sigma_jL_{L^2+j-1}^{(1)}(u).         \tag{9.5}
\]
Since
\[
 p_n(\eps)=\int_0^\infty e^{-\eps u}L_{n-1}^{(1)}(u)\,du,
\]
finite-dimensional \(\ell^1\)--\(\ell^\infty\) duality gives
\[
\boxed{
 S_L=\lim_{\eps\downarrow0}\sup_{\sigma\in[-1,1]^L}
 \left\{
 \sum_{j=0}^{L-1}\sigma_jA_{L^2+j}
 +\int_0^\infty P_{\sigma,L}(u)\,d\nu_\eps(u)
 \right\}.}                                             \tag{9.6}
\]
The interchange is legitimate because the block is finite: the difference
between the two suprema is at most
\(\sum_{j<L}|\lambda_{L^2+j,\eps}-\lambda_{L^2+j}|\).
Consequently, the direct ordinary-prime closure corresponding to (8.9) is
\[
 \log\!\left[1+\lim_{\eps\downarrow0}
 \sup_{\sigma\in[-1,1]^L}
 \left\{
 \sum_{j<L}\sigma_jA_{L^2+j}
 +\int P_{\sigma,L}\,d\nu_\eps
 \right\}\right]=o(L^2).                               \tag{9.7}
\]
By Theorem~\ref{thm:block-radius}, (9.7) is exactly equivalent to RH.
It is not inferred by placing separate moduli on the continuum and prime
parts of (9.4).

\subsection{Lossless prefix compression}

Put \(N=L^2\) and
\[
 T_{L,j}=\sum_{r=0}^{j}\lambda_{N+r},\qquad
 M_L=\max_{0\le j<L}|T_{L,j}|.                          \tag{9.8}
\]

\begin{proposition}[Prefix compression]
\label{prop:prefix-compression}
For every \(L\ge1\),
\[
 \boxed{M_L\le S_L\le(2L-1)M_L.}                       \tag{9.9}
\]
Consequently
\[
 \boxed{\RH\Longleftrightarrow\log(1+M_L)=o(L^2).}     \tag{9.10}
\]
\end{proposition}

\begin{proof}
The first inequality is the triangle inequality. Conversely, if
\(x_r=\lambda_{N+r}\), then \(x_0=T_{L,0}\) and
\(x_r=T_{L,r}-T_{L,r-1}\) for \(r\ge1\). Hence
\[
 \sum_{r<L}|x_r|
 \le |T_{L,0}|+\sum_{r=1}^{L-1}
 (|T_{L,r}|+|T_{L,r-1}|)
 \le(2L-1)M_L.
\]
Since \(\log(2L-1)=o(L^2)\), (9.10) follows from
Theorem~\ref{thm:block-radius}.
\end{proof}

This elementary compression has a useful exact arithmetic consequence.
Define
\[
 \Delta(u)=\gamma-u+\sum_{m\le e^u}\frac{\Lambda(m)}m,
 \qquad
 H_{N,j}(u)=L_{N+j-2}^{(3)}(u)-L_{N-3}^{(3)}(u).        \tag{9.11}
\]
Abel summation of the coupled measure (9.4), followed by
\[
 \sum_{r=0}^{j}L_{N+r-1}^{(1)}
 =L_{N+j-1}^{(2)}-L_{N-2}^{(2)},                        \tag{9.12}
\]
gives
\[
\begin{split}
 T_{L,j}={}&\sum_{r=0}^{j}A_{N+r}
 +\gamma\left\{
 \binom{N+j+1}{2}-\binom N2\right\}\\
 &-\int_0^\infty H_{N,j}(u)\Delta(u)\,du.
\end{split}                                               \tag{9.13}
\]
The two explicit terms are \(\exp\{o(N)\}\), uniformly for
\(j<L=\sqrt N\). Therefore the remaining ordinary-prime estimate is
precisely
\[
 \boxed{
 \log\!\left(1+\max_{0\le j<L}\left|
 \int_0^\infty
 \bigl(L_{L^2+j-2}^{(3)}-L_{L^2-3}^{(3)}\bigr)(u)
 \left(\gamma-u+\sum_{m\le e^u}\frac{\Lambda(m)}m\right)du
 \right|\right)=o(L^2).}                                \tag{9.14}
\]
By (9.10)--(9.13), (9.14) is equivalent to RH.  Its advantage over
(9.7) is combinatorial rather than logical: only \(L\) two-boundary
kernels remain.  No estimate for (9.14) is asserted here.

\begin{problem}[Arithmetic dissipativity]
\label{prob:closure}
Prove one of the following directly for the ordinary von Mangoldt weights:
\begin{enumerate}
\item for one fixed \(L\), inequality \((9.1)\) holds for all sufficiently
      large \(N\);
\item there exist consecutive intervals \(I_k\), with
      \(|I_k|\to\infty\), such that \(\mathcal S(I_k)=0\);
\item for one fixed \(L\), the adjacent-bank energy satisfies
      \(\mathcal E_L(N)\le e^{o(N)}\).
\end{enumerate}
Each assertion proves RH.
\end{problem}

Indeed, item 1 closes by Theorem~\ref{thm:window-positive}, and item 3 by
Corollary~\ref{cor:energy}. For item 2, if RH were false then
Corollary~\ref{cor:two-sided} would provide a syndetic set \(D_+\) on
which \(\lambda_n\ge cR^n\). Eventually \(cR^n>e^{\sqrt n}\), and every
sufficiently long interval would meet \(D_+\); hence its supply
\(\mathcal S(I)\) would be strictly positive. This contradicts the
existence of zero-supply intervals of unbounded length.

The three targets encode unstable-mode exclusion with different
nonlinearities: a half-space, a positive-part supply, and an energy upper
bound. The first is closest to classical dissipativity; the second is the
weakest pointwise block formulation; the third is phase blind.

\section{Status table}

The results separate observability from stability.

\begin{center}
\begin{tabular}{p{0.30\textwidth}p{0.17\textwidth}p{0.36\textwidth}}
\toprule
Statement & Status & Meaning\\
\midrule
Off-line quartet has exponential envelope & proved & unstable mode exists\\
Other zeros cannot cancel it forever & proved & global observability\\
Moving averages preserve it & proved & FIR detector has no off-unit zero\\
Adjacent windows have no common blind spot & proved & quantitative observability\\
Vanishing-trim and geometric-rate criteria & proved & robust nonlinear detection\\
Exact absolute-block spectral radius & proved & block rate equals largest Cayley modulus\\
Prime--Laguerre window identity & proved & exact arithmetic realization\\
Signed average implies a good pointwise block & false & rational quartet witness\\
Ordinary-prime FIR output is subexponential & open & stability estimate\\
Ordinary-prime FIR output is eventually nonnegative & open & one-sided dissipativity\\
\bottomrule
\end{tabular}
\end{center}

The fact that all computed zeros are critical validates the marginally
stable response over the verified spectral range. It does not prove the
last two rows. Similarly, the exact decomposition into primes, pole and
Gamma factor identifies every component of the plant, but a decomposition
does not provide the signed upper estimate needed to establish stability.

\section{Reproducibility}

The companion script
\begin{center}
\path{03-research/phase-105-a1-geometry-and-deep-limit/tools/}\
\path{fir_window_bank_check.py}
\end{center}
uses exact rational arithmetic for the Laguerre degree-sum identity and
checks the FIR identities, the adjacent-bank inequality, survival of an
off-unit mode through a cascade, and the rational signed-average witness.
It is run from the repository root with
\begin{center}
\texttt{python3}\
\path{03-research/phase-105-a1-geometry-and-deep-limit/tools/fir_window_bank_check.py}.
\end{center}
These computations audit algebraic identities; the analytic proofs are the
arguments in the preceding sections.

\section{Conclusion}

The Cayley transform makes the signals-and-systems content of the Li
criterion exact. Critical zeros are unit-circle modes. Off-line zeros are
unstable modes. The Li coefficients are the sampled output. Moving windows
are FIR filters. Repeated averaging is a filter cascade. A pair of adjacent
windows is an observable filter bank with the quantitative lower bound
\((5.2)\).

For every fixed \(L\),
\[
 \RH\iff
 \limsup_N\left|\sum_{j=0}^{L-1}\lambda_{N+j}\right|^{1/N}
 \le1,
\]
and also
\[
 \RH\iff
 \sum_{j=0}^{L-1}\lambda_{N+j}\ge0
 \quad\hbox{for every sufficiently large }N.
\]
The corresponding prime-side output uses one difference of two Laguerre
polynomials of parameter \(2\).

The remaining issue is not whether the filter detects an off-line zero; it
does. It is whether the literal arithmetic plant satisfies a subexponential
upper bound or a one-sided supply inequality. Repeated linear averaging
alone cannot prove that bound, and a quadratic detector still requires an
upper estimate after exact prime--pole cancellation.
Problem~\ref{prob:closure} records the closure in three control-theoretic
forms. Solving any one of them would complete a proof of RH.

\begin{thebibliography}{99}

\bibitem{Titchmarsh}
E.~C.~Titchmarsh,
\emph{The Theory of the Riemann Zeta-Function}, 2nd ed.,
revised by D.~R.~Heath-Brown, Oxford University Press, 1986.

\bibitem{Keiper}
J.~B.~Keiper,
\emph{Power series expansions of Riemann's \(\xi\) function},
Math. Comp. \textbf{58} (1992), 765--773.

\bibitem{Li}
X.-J.~Li,
\emph{The positivity of a sequence of numbers and the Riemann hypothesis},
J. Number Theory \textbf{65} (1997), 325--333.

\bibitem{BombieriLagarias}
E.~Bombieri and J.~C.~Lagarias,
\emph{Complements to Li's criterion for the Riemann hypothesis},
J. Number Theory \textbf{77} (1999), 274--287.

\bibitem{Lagarias}
J.~C.~Lagarias,
\emph{Li coefficients for automorphic \(L\)-functions},
Ann. Inst. Fourier (Grenoble) \textbf{57} (2007), 1689--1740.

\bibitem{Coffey}
M.~W.~Coffey,
\emph{Toward verification of the Riemann hypothesis: application of the
Li criterion},
Math. Phys. Anal. Geom. \textbf{8} (2005), 211--255.

\bibitem{Szego}
G.~Szeg\H{o},
\emph{Orthogonal Polynomials}, 4th ed.,
American Mathematical Society, 1975.

\bibitem{Oppenheim}
A.~V.~Oppenheim and R.~W.~Schafer,
\emph{Discrete-Time Signal Processing}, 3rd ed.,
Pearson, 2010.

\bibitem{Kailath}
T.~Kailath,
\emph{Linear Systems},
Prentice--Hall, 1980.

\bibitem{ZhouDoyleGlover}
K.~Zhou, J.~C.~Doyle, and K.~Glover,
\emph{Robust and Optimal Control},
Prentice--Hall, 1996.

\end{thebibliography}

\end{document}
